% GENERATED FILE. Edit canonical lessons, then run npm run generate:books.
% canonical-source-hash: fnv1a64:1e1b5d905970d5d1
% canonical-lessons: SA-C107-gudah, SA-C107-maricam, SA-C107-haridra, SA-C107-ardrakam, SA-C107-lasunam

\chapter{Jaggery, Pepper, Turmeric, Ginger, Garlic}
\label{ch:sa-jaggery-pepper-turmeric-ginger-garlic}
\hlchaptermodality{107}

\begin{chapteropening}
\textbf{By the end of this chapter:} \emph{I can say jaggery, pepper, turmeric, ginger and garlic.}

Complete the last lesson of chapter 107: i can say jaggery, pepper, turmeric, ginger and garlic.
\end{chapteropening}

\section[\texorpdfstring{\sk{गुडः} (guḍaḥ)}{गुडः (guḍaḥ)}]{\sk{गुडः} (guḍaḥ) — jaggery}

\label{lesson:SA-C107-gudah}



\begin{quote}
Before the new one: say the Sanskrit for mung beans, then the Sanskrit for dal.
\end{quote}

\subsection*{You'll want to know: \sk{गुडः}}
\textbf{\sk{गुडः}} — \emph{guḍaḥ} — \textquotedblleft{}jaggery\textquotedblright{}.

\begin{grammarlens}[title={What you've built}]
One more word for this chapter.
\end{grammarlens}

\subsection*{Your turn}
\begin{itemize}
\raggedright
  \item \emph{Say it:} \emph{guḍaḥ}
  \item \emph{Say it:} \emph{guḍaḥ}, once more
  \item \emph{Say it:} say \emph{guḍaḥ}
  \item \emph{Recall:} say the Sanskrit for mung beans, then the Sanskrit for dal, then say \emph{guḍaḥ} again
\end{itemize}

\begin{tcolorbox}[breakable,colback=teal!4,colframe=teal!35!black,title={Before you move on}]
What does \sk{गुडः} mean? (\textquotedblleft{}Jaggery\textquotedblright{}.) Say it once more.
\end{tcolorbox}
\section[\texorpdfstring{\sk{मरीचम्} (marīcam)}{मरीचम् (marīcam)}]{\sk{मरीचम्} (marīcam) — pepper}

\label{lesson:SA-C107-maricam}



\begin{quote}
Before the new one: say the Sanskrit for dal, then the Sanskrit for jaggery.
\end{quote}

\subsection*{You'll want to know: \sk{मरीचम्}}
\textbf{\sk{मरीचम्}} — \emph{marīcam} — \textquotedblleft{}pepper\textquotedblright{}.

\begin{grammarlens}[title={What you've built}]
One more word for this chapter.
\end{grammarlens}

\subsection*{Your turn}
\begin{itemize}
\raggedright
  \item \emph{Say it:} \emph{marīcam}
  \item \emph{Say it:} \emph{marīcam}, once more
  \item \emph{Say it:} say \emph{marīcam}
  \item \emph{Recall:} say the Sanskrit for dal, then the Sanskrit for jaggery, then say \emph{marīcam} again
\end{itemize}

\begin{tcolorbox}[breakable,colback=teal!4,colframe=teal!35!black,title={Before you move on}]
What does \sk{मरीचम्} mean? (\textquotedblleft{}Pepper\textquotedblright{}.) Say it once more.
\end{tcolorbox}
\section[\texorpdfstring{\sk{हरिद्रा} (haridrā)}{हरिद्रा (haridrā)}]{\sk{हरिद्रा} (haridrā) — turmeric}

\label{lesson:SA-C107-haridra}



\begin{quote}
Before the new one: say the Sanskrit for jaggery, then the Sanskrit for pepper.
\end{quote}

\subsection*{You'll want to know: \sk{हरिद्रा}}
\textbf{\sk{हरिद्रा}} — \emph{haridrā} — \textquotedblleft{}turmeric\textquotedblright{}.

\begin{grammarlens}[title={What you've built}]
One more word for this chapter.
\end{grammarlens}

\subsection*{Your turn}
\begin{itemize}
\raggedright
  \item \emph{Say it:} \emph{haridrā}
  \item \emph{Say it:} \emph{haridrā}, once more
  \item \emph{Say it:} say \emph{haridrā}
  \item \emph{Recall:} say the Sanskrit for jaggery, then the Sanskrit for pepper, then say \emph{haridrā} again
\end{itemize}

\begin{tcolorbox}[breakable,colback=teal!4,colframe=teal!35!black,title={Before you move on}]
What does \sk{हरिद्रा} mean? (\textquotedblleft{}Turmeric\textquotedblright{}.) Say it once more.
\end{tcolorbox}
\section[\texorpdfstring{\sk{आर्द्रकम्} (ārdrakam)}{आर्द्रकम् (ārdrakam)}]{\sk{आर्द्रकम्} (ārdrakam) — ginger}

\label{lesson:SA-C107-ardrakam}



\begin{quote}
Before the new one: say the Sanskrit for pepper, then the Sanskrit for turmeric.
\end{quote}

\subsection*{You'll want to know: \sk{आर्द्रकम्}}
\textbf{\sk{आर्द्रकम्}} — \emph{ārdrakam} — \textquotedblleft{}ginger\textquotedblright{}.

\begin{grammarlens}[title={What you've built}]
One more word for this chapter.
\end{grammarlens}

\subsection*{Your turn}
\begin{itemize}
\raggedright
  \item \emph{Say it:} \emph{ārdrakam}
  \item \emph{Say it:} \emph{ārdrakam}, once more
  \item \emph{Say it:} say \emph{ārdrakam}
  \item \emph{Recall:} say the Sanskrit for pepper, then the Sanskrit for turmeric, then say \emph{ārdrakam} again
\end{itemize}

\begin{tcolorbox}[breakable,colback=teal!4,colframe=teal!35!black,title={Before you move on}]
What does \sk{आर्द्रकम्} mean? (\textquotedblleft{}Ginger\textquotedblright{}.) Say it once more.
\end{tcolorbox}
\section[\texorpdfstring{\sk{लशुनम्} (laśunam)}{लशुनम् (laśunam)}]{\sk{लशुनम्} (laśunam) — garlic}

\label{lesson:SA-C107-lasunam}



\begin{quote}
Before the new one: say the Sanskrit for turmeric, then the Sanskrit for ginger.
\end{quote}

\subsection*{You'll want to know: \sk{लशुनम्}}
\textbf{\sk{लशुनम्}} — \emph{laśunam} — \textquotedblleft{}garlic\textquotedblright{}.

\begin{grammarlens}[title={What you've built}]
One more word for this chapter.
\end{grammarlens}

\subsection*{Your turn}
\begin{itemize}
\raggedright
  \item \emph{Say it:} \emph{laśunam}
  \item \emph{Say it:} \emph{laśunam}, once more
  \item \emph{Say it:} say \emph{laśunam}
  \item \emph{Recall:} say the Sanskrit for turmeric, then the Sanskrit for ginger, then say \emph{laśunam} again
\end{itemize}

\begin{tcolorbox}[breakable,colback=teal!4,colframe=teal!35!black,title={Before you move on}]
What does \sk{लशुनम्} mean? (\textquotedblleft{}Garlic\textquotedblright{}.) Say it once more.
\end{tcolorbox}
